\section{Terminal base change and the solution regular image}
\label{sec:terminalization}

This section completes the proof over the total parameter object.  Coproducts retain all admissible data, regular images express existential projection, and pullbacks perform the terminal base change.

Define the admissible parameter object
\begin{equation}\label{eq:terminal-parameter-object}
\mathfrak P
=
\left\{(\eta,\sigma,\rho):
0<\eta<\eta_0,\quad
0<\sigma<1,\quad
0<\rho<\kappa\sigma\rho_P
\right\}.
\end{equation}
It is nonempty.  For $a=(\eta,\sigma,\rho)\in\mathfrak P$, write
\[
R_a=R_\rho,
\qquad
\mathcal B_{\mathrm{phys}}(a)
=\mathcal B_{\mathrm{phys}}(\eta,\sigma,\rho),
\]
and define the physically admissible scalar-start object
\begin{equation}\label{eq:terminal-good-start-object}
\mathfrak G_a
=
\left\{(B,\xi):
(B,\xi)\in\mathfrak G_{\mathrm{sc}}(\eta,\sigma,\rho),\quad
B\in\mathcal B_{\mathrm{phys}}(\eta,\sigma,\rho)
\right\}.
\end{equation}
Theorem~\ref{thm:scalar-terminality} says that the $B$-image of $\mathfrak G_{\mathrm{sc}}(\eta,\sigma,\rho)$ is cofinite, whereas \eqref{eq:global-start-cutoff} contains a final segment.  Their intersection is therefore nonempty, so every $\mathfrak G_a$ is nonempty.  Hence the coproduct
\begin{equation}\label{eq:terminal-total-start-object}
\mathfrak G
=
\coprod_{a\in\mathfrak P}\mathfrak G_a
\end{equation}
projects by a regular epimorphism to $\mathfrak P$.

For $g=(a,B,\xi)\in\mathfrak G$, retain the entire nonempty scalar-threshold object $\mathfrak K_\xi$ from \eqref{eq:scalar-threshold-object}.  Form
\begin{equation}\label{eq:terminal-threshold-object}
\mathfrak K
=
\coprod_{g=(a,B,\xi)\in\mathfrak G}\mathfrak K_\xi.
\end{equation}
The projection $\mathfrak K\twoheadrightarrow\mathfrak G$ is a regular epimorphism.

Above the total threshold object form its universal tail
\begin{equation}\label{eq:terminal-threshold-tail}
\mathfrak T_{\mathfrak K}
=
\{(g,k_0,k):(g,k_0)\in\mathfrak K,\ k\ge k_0\}.
\end{equation}
The projection $\mathfrak T_{\mathfrak K}\twoheadrightarrow\mathfrak K$ is a regular epimorphism; its fibre over $(g,k_0)$ contains the canonical diagonal count $k=k_0$.  Now retain all cutoff witnesses over this tail:
\begin{equation}\label{eq:universal-cutoff-incidence}
\mathfrak X_*
=
\{(g,k_0,k,X):(g,k_0,k)\in\mathfrak T_{\mathfrak K},\ (k,X)\in\mathfrak X_\xi\},
\end{equation}
where $g=(a,B,\xi)$.  By the defining property of $k_0\in\mathfrak K_\xi$, the projection
\begin{equation}\label{eq:cutoff-threshold-regular-epi}
r_*:\mathfrak X_*\twoheadrightarrow\mathfrak T_{\mathfrak K},
\qquad
(g,k_0,k,X)\longmapsto(g,k_0,k)
\end{equation}
is a regular epimorphism.  The structural existence maps therefore paste to the regular-epimorphism chain
\begin{equation}\label{eq:terminal-existence-chain}
\mathfrak X_*
\mathrel{\twoheadrightarrow}
\mathfrak T_{\mathfrak K}
\mathrel{\twoheadrightarrow}
\mathfrak K
\mathrel{\twoheadrightarrow}
\mathfrak G
\mathrel{\twoheadrightarrow}
\mathfrak P
\mathrel{\twoheadrightarrow}
\{*\}.
\end{equation}
Thus every parameter object in \eqref{eq:terminal-existence-chain} is nonempty, and the existential structure is carried by the regular epimorphisms themselves.

Let
\[
c:\mathfrak T_{\mathfrak K}\longrightarrow\N_{>0},
\qquad
(g,k_0,k)\longmapsto k,
\]
and let
\begin{equation}\label{eq:terminal-count-image}
J=\operatorname{im}(c)\hookrightarrow\N_{>0}
\end{equation}
be its regular image.  By construction every threshold point contributes its entire upper tail, so $J$ is upward closed; by \eqref{eq:terminal-existence-chain} it is nonempty.  Hence $J$ is a final segment and in particular is cofinite.  The image factorization of $c$, pasted with \eqref{eq:cutoff-threshold-regular-epi}, gives
\begin{equation}\label{eq:cutoff-count-regular-epi}
\bar p_*:\mathfrak X_*\twoheadrightarrow J.
\end{equation}

An element $x=(g,k_0,k,X)\in\mathfrak X_*$ contains $g=(a,B,\xi)$ and satisfies
\[
\gamma_X^\xi=0,
\qquad
k\in I_X^\xi.
\]
Consequently
\[
(0,k)\in
T_X^\xi
=
\operatorname{Fib}_{E_X}(0)\times I_X^\xi.
\]
Moreover $B\in\mathcal B_{\mathrm{phys}}(a)$, so Proposition~\ref{prop:finite-realization} supplies the universal covering correspondence
\[
R_{X,\xi}^{\mathrm{univ}}:\{0_M\}\rightsquigarrow T_X^\xi.
\]
Form the total edge and target objects over the whole cutoff incidence:
\[
\mathfrak R_*
=
\coprod_{x=(g,k_0,k,X)\in\mathfrak X_*}
E_{R_{X,\xi}^{\mathrm{univ}}},
\qquad
\mathfrak T_*
=
\coprod_{x=(g,k_0,k,X)\in\mathfrak X_*}T_X^\xi.
\]
The target maps assemble to a regular epimorphism
\[
\tau_*:\mathfrak R_*\twoheadrightarrow\mathfrak T_*.
\]
There is a canonical point-of-the-fibre map
\[
j_*:\mathfrak X_*\longrightarrow\mathfrak T_*,
\qquad
(g,k_0,k,X)\longmapsto(g,k_0,k,X;(0,k)).
\]
Define the total terminal edge object by the cartesian square
\begin{equation}\label{eq:terminal-edge-pullback}
\begin{CD}
\mathfrak E_* @>>> \mathfrak R_*\\
@V{e_*}VV @VV{\tau_*}V\\
\mathfrak X_* @>{j_*}>> \mathfrak T_*.
\end{CD}
\end{equation}
Since regular epimorphisms are stable under pullback,
\[
e_*:\mathfrak E_*\twoheadrightarrow\mathfrak X_*
\]
is a regular epimorphism.  Pasting with \eqref{eq:cutoff-count-regular-epi} gives
\begin{equation}\label{eq:terminal-witness-regular-epi}
\bar p_*e_*:\mathfrak E_*\twoheadrightarrow J.
\end{equation}
This regular epimorphism is the existential content of the proof.

Let
\[
q_{\Z}:\Q\longrightarrow\Q/\Z
\]
be the quotient map.  By definition of $\overline W$, the square
\begin{equation}\label{eq:value-residue-square}
\begin{CD}
\mathcal C @>{W}>> \Q\\
@V{\overline W}VV @VV{q_{\Z}}V\\
\Q/\Z @= \Q/\Z
\end{CD}
\end{equation}
commutes.  The total edge object has a label map
\[
\ell_*:\mathfrak E_*\longrightarrow\mathcal C.
\]
For every $e\in\mathfrak E_*$ lying over a cutoff incidence of count $k$, the pullback \eqref{eq:terminal-edge-pullback} and the conservation square \eqref{eq:edge-conservation-square} give
\[
\overline W(\ell_*(e))=0,
\qquad
\nu(\ell_*(e))=k.
\]
Proposition~\ref{prop:finite-realization} gives $W(\ell_*(e))<2$.  Since $k>0$, the label is nonempty and its reciprocal value is positive.  The square \eqref{eq:value-residue-square} therefore gives
\[
0<W(\ell_*(e))<2,
\qquad
W(\ell_*(e))\in\ker q_{\Z}=\Z,
\]
and hence $W(\ell_*(e))=1$.

Thus the map
\[
e\longmapsto\bigl(\bar p_*e_*(e),\ell_*(e)\bigr)
\]
factors through the solution object $\mathscr S$ of \eqref{eq:solution-object}.  We obtain the commuting triangle
\begin{equation}\label{eq:solution-factorization}
\begin{CD}
\mathfrak E_* @>>> \mathscr S\\
@V{\bar p_*e_*}VV @VV{\pi_{\mathrm{sol}}}V\\
J @>>> \N_{>0},
\end{CD}
\end{equation}
where the lower arrow is the regular-image inclusion.  The left vertical map is a regular epimorphism, so
\[
J\subseteq\operatorname{im}(\pi_{\mathrm{sol}}).
\]
Since $J$ is cofinite, the regular image of $\pi_{\mathrm{sol}}$ is cofinite.  This is Theorem~\ref{thm:main}.
